\documentclass[11pt,a4paper]{article}

% The document is written with XeTeX/Tectonic so that the Greek explanation
% and the mathematical notation can coexist without transliteration.
\usepackage{fontspec}
\setmainfont{Times New Roman}
\setsansfont{Arial}
\usepackage[a4paper,margin=24mm,headheight=24pt]{geometry}
\usepackage{amsmath}
\usepackage{booktabs}
\usepackage{enumitem}
\usepackage{xcolor}
\usepackage{fancyhdr}
\usepackage{hyperref}
\usepackage{titlesec}

\definecolor{navy}{HTML}{102A43}
\definecolor{bodyblue}{HTML}{243B53}
\definecolor{muted}{HTML}{486581}
\definecolor{ruleblue}{HTML}{D9E2EC}
\definecolor{boxblue}{HTML}{F0F4F8}
\definecolor{codeblue}{HTML}{16324F}

\hypersetup{
  colorlinks=true,
  linkcolor=navy,
  urlcolor=navy,
  pdftitle={Forecast.py: Step-by-step mathematical walkthrough},
  pdfauthor={Project Optimus}
}

\pagestyle{fancy}
\fancyhf{}
\lhead{\textcolor{muted}{IAM Workload Capacity Model}}
\rhead{\textcolor{muted}{forecast.py walkthrough}}
\cfoot{\textcolor{muted}{\thepage}}
\renewcommand{\headrulewidth}{0.4pt}
\renewcommand{\headrule}{\hbox to\headwidth{\color{ruleblue}\leaders\hrule height \headrulewidth\hfill}}

\titleformat{\section}{\Large\bfseries\color{navy}}{\thesection.}{0.6em}{}
\titleformat{\subsection}{\large\bfseries\color{bodyblue}}{\thesubsection}{0.6em}{}
\setlist[itemize]{leftmargin=8mm,itemsep=2pt,topsep=3pt}
\setlist[enumerate]{leftmargin=8mm,itemsep=3pt,topsep=3pt}
\setlength{\parindent}{0pt}
\setlength{\parskip}{6pt}
\color{bodyblue}

\newcommand{\equationnote}[1]{
  \begin{center}
    \colorbox{boxblue}{\parbox{0.92\linewidth}{\begin{equation*}#1\end{equation*}}}
  \end{center}
}

\newcommand{\code}[1]{\texttt{#1}}

\begin{document}

\begin{titlepage}
  \centering
  \vspace*{28mm}
  {\Huge\bfseries\color{navy} Από τον κώδικα στις φόρμουλες\\[4mm]
  \Large Το \code{forecast.py} βήμα προς βήμα\par}
  \vspace{10mm}
  {\large\color{muted} IAM request-demand και workload forecasting\par}
  \vspace{18mm}
  \fcolorbox{ruleblue}{boxblue}{\begin{tabular}{ccc}
    \textbf{Project Optimus} & \textbf{Version 0.1} & \textbf{Model walkthrough}
  \end{tabular}}
  \vfill
  {\color{muted} Το έγγραφο εξηγεί τι κάνει η τρέχουσα υλοποίηση,\
  όχι τι θα έκανε μια μελλοντική production έκδοση.\par}
\end{titlepage}

\tableofcontents
\newpage

\section{Τι προσπαθεί να προβλέψει ο κώδικας;}

Ο κώδικας δεν προβλέπει απευθείας ένα μόνο νούμερο ``συνολικού κόστους''.
Κατασκευάζει πρώτα μια \textbf{κατανομή πιθανών μελλοντικών workloads}. Αυτό
είναι σημαντικό: ένα capacity plan χρειάζεται όχι μόνο τη μέση ζήτηση, αλλά
και μια εικόνα για την αβεβαιότητα και τις υψηλές ουρές της ζήτησης.

Για ένα μελλοντικό time bucket $t$ (π.χ. μία ημέρα), η συνολική ζήτηση
workload είναι:

\equationnote{W_t = \sum_{i=1}^{N_t} w_{i,t}}

όπου:
\begin{itemize}
  \item $N_t$ είναι ο αριθμός των requests στο bucket $t$,
  \item $w_{i,t}$ είναι το workload του request $i$ σε capacity-unit time,
  \item $W_t$ είναι το συνολικό request workload του bucket.
\end{itemize}

Το \code{forecast.py} παράγει πολλά πιθανά $W_t$ μέσω simulation. Η μέση
τιμή ή το P95 προκύπτουν αργότερα από αυτά τα simulated totals.

\section{Τα τρία συστατικά του compound forecast}

Το ``compound'' σημαίνει ότι το συνολικό workload αποτελείται από τυχαίο
αριθμό requests, όπου κάθε request έχει τυχαίο outcome και τυχαίο βάρος.
Ο κώδικας συνδυάζει τρία μοντέλα:

\begin{enumerate}
  \item \textbf{Request-count model:} προβλέπει πόσα requests θα έρθουν,
        δηλαδή το $N_t$.
  \item \textbf{Rejection-rate model:} προβλέπει ποιο ποσοστό θα απορριφθεί,
        δηλαδή την πιθανότητα $p_{R,t}$.
  \item \textbf{Conditional workload distributions:} δίνουν το πιθανό
        workload ενός approved ή rejected request.
\end{enumerate}

Η τελική προσομοίωση είναι επομένως:

\equationnote{
  N_t \longrightarrow O_{1,t},\ldots,O_{N_t,t}
  \longrightarrow w_{1,t},\ldots,w_{N_t,t}
  \longrightarrow W_t = \sum_i w_{i,t}
}

Τα βέλη σημαίνουν ``κάνε ένα random draw από το προηγούμενο μοντέλο''. Δεν
κάνουμε ένα αυθαίρετο άλμα από το forecast του $N_t$ στο workload: παράγουμε
ρητά τα request weights και τα αθροίζουμε.

\section{Πριν από το \code{forecast.py}: το workload ενός request}

Το \code{forecast.py} δεν υπολογίζει από την αρχή το κόστος κάθε ιστορικού
request. Παραλαμβάνει τις κατανομές που έχουν δημιουργηθεί στο
\code{workload.py}. Η βάση αυτών των κατανομών είναι το workload contract του
\code{models.py}.

\subsection{Backend tasks}

Αν το request $i$ δημιουργεί backend tasks $j=1,\ldots,J_i$, κάθε task έχει:

\begin{itemize}
  \item $\tau_{ij}$: active execution time σε δευτερόλεπτα,
  \item $r_{ij}$: capacity units που καταναλώνει όσο εκτελείται.
\end{itemize}

Το workload του task είναι το γινόμενο χρόνου επί πόρων:

\equationnote{w_{ij}^{task} = \tau_{ij} r_{ij}}

Αν δύο tasks τρέχουν παράλληλα, τα αθροίζουμε. Ο λόγος είναι ότι μπορεί να
μην αυξάνεται το wall-clock duration του request, αλλά δεσμεύονται ταυτόχρονα
δύο capacity-consuming executions.

Ο κώδικας υπολογίζει επίσης προαιρετικό orchestration/workflow overhead:

\equationnote{w_i^{wf} = T_i^{wf} r_i^{wf}}

όπου $T_i^{wf}$ είναι ο ενεργός χρόνος του workflow engine και $r_i^{wf}$ η
ένταση κατανάλωσης capacity. Ο χρόνος αναμονής για manual approval δεν
μπαίνει εδώ, εκτός αν η συγκεκριμένη αναμονή κρατά πράγματι δεσμευμένο
capacity resource.

Άρα το backend μέρος είναι:

\equationnote{
  w_i^{backend} = \sum_{j=1}^{J_i} \tau_{ij}r_{ij} + T_i^{wf}r_i^{wf}
}

\subsection{Approval outcome και τελικό request weight}

Το $I_i^{A}$ είναι indicator μεταβλητή: παίρνει την τιμή $1$ όταν το request
εγκρίνεται και $0$ όταν απορρίπτεται.

\equationnote{
  I_i^{A} =
  \begin{cases}
    1, & \text{αν το request εγκρίθηκε},\\
    0, & \text{αν το request απορρίφθηκε}.
  \end{cases}
}

Αν υπάρχει κόστος μέχρι να ληφθεί η απόφαση, το ονομάζουμε
$w_i^{pre}$. Το συνολικό observed weight είναι:

\equationnote{w_i = w_i^{pre} + I_i^{A}w_i^{backend}}

Έτσι ένα rejected request διατηρεί το pre-decision κόστος, αλλά δεν παίρνει
backend calculation cost. Αυτή ακριβώς η διάκριση συνδέει αργότερα το
rejection forecast με το capacity forecast.

\section{Βήμα 1: fitting του request count}

Η κλάση \code{NegativeBinomialModel} ξεκινά από ιστορικούς αριθμούς requests
ανά time bucket. Για παρατηρήσεις $n_1,\ldots,n_T$ υπολογίζει πρώτα τον μέσο
αριθμό requests:

\equationnote{\hat{\mu} = \frac{1}{T}\sum_{t=1}^{T} n_t}

και τη διακύμανση με τον ίδιο population-style παρονομαστή που χρησιμοποιεί
ο κώδικας:

\equationnote{\hat{v} = \frac{1}{T}\sum_{t=1}^{T}(n_t-\hat{\mu})^2}

\subsection{Γιατί Negative Binomial;}

Για Negative Binomial με mean $\mu$ και dispersion $\phi$ χρησιμοποιούμε την
παραμετροποίηση:

\equationnote{\operatorname{Var}(N_t) = \mu + \phi\mu^2}

Το πρώτο μέρος, $\mu$, είναι η διακύμανση του Poisson. Το δεύτερο μέρος,
$\phi\mu^2$, είναι επιπλέον μεταβλητότητα. Επομένως το $\phi$ επιτρέπει στα
request counts να είναι πιο ασταθή από ένα Poisson process.

Λύνουμε την παραπάνω σχέση ως προς $\phi$:

\begin{align*}
  v &= \mu + \phi\mu^2,\\
  v-\mu &= \phi\mu^2,\\
  \hat{\phi} &= \frac{\hat{v}-\hat{\mu}}{\hat{\mu}^2}.
\end{align*}

Αυτό είναι το method-of-moments fitting που εφαρμόζει το
\code{NegativeBinomialModel.fit}. Αν $\hat{v}\leq\hat{\mu}$, η εκτίμηση θα
ήταν αρνητική, κάτι που δεν επιτρέπεται ως dispersion. Ο κώδικας την κόβει
στο μηδέν:

\equationnote{\hat{\phi} = \max\left(0,\frac{\hat{v}-\hat{\mu}}{\hat{\mu}^2}\right)}

Με $\phi=0$ το μοντέλο γίνεται το Poisson limit. Άρα η υλοποίηση δεν
επιβάλλει υπερδιασπορά όταν τα δεδομένα δεν την υποστηρίζουν.

\textbf{Σημαντικό:} αυτό είναι αρχικό baseline fitting, όχι πλήρης maximum
likelihood optimization. Το πλεονέκτημα είναι ότι η λογική μένει ορατή και
χωρίς numerical dependency. Σε επόμενο στάδιο μπορεί να αντικατασταθεί ο
εκτιμητής χωρίς να αλλάξει το interface του compound forecast.

\section{Βήμα 2: από πού προκύπτει το $s_{t+1}$;}

Αυτό είναι το σημείο που χρειάζεται τη μεγαλύτερη διευκρίνιση. Η σχέση

\equationnote{s_{t+1}=\omega+\text{persistence}\,s_t+\text{learning\_rate}\,\text{score}_t}

δεν προκύπτει αυτόματα μόνο και μόνο επειδή επιλέξαμε Negative Binomial.
Είναι ένας \textbf{επιλεγμένος δυναμικός μηχανισμός online correction},
βασισμένος στη λογική των score-driven models. Ο σκοπός του είναι να
ανανεώνει το forecast όταν κλείνει ένα νέο time bucket, χωρίς να ξανακάνει
fit όλο το ιστορικό από την αρχή.

\subsection{Το state και η σύνδεσή του με το mean}

Το state $s_t$ δεν είναι workload ούτε αριθμός requests. Είναι latent level
στη λογαριθμική κλίμακα. Η σύνδεση είναι:

\equationnote{s_t = \log(\mu_t), \qquad \mu_t = \exp(s_t)}

Χρησιμοποιούμε το log διότι το $\mu_t$ πρέπει να είναι μη αρνητικό. Όποια
πραγματική τιμή και αν έχει το $s_t$, η εκθετική μετατροπή παράγει θετικό
forecast mean. Στον κώδικα το αρχικό state μετά το fit είναι:

\equationnote{s_1 = \log(\hat{\mu})}

και το \code{forecast()} επιστρέφει $\exp(s_t)$.

\subsection{Το score που χρησιμοποιεί η τρέχουσα υλοποίηση}

Με forecast mean $\mu_t$ και observed count $N_t$, το raw forecast error είναι
$N_t-\mu_t$. Ο κώδικας το κανονικοποιεί ως:

\equationnote{
  \text{score}_t =
  \frac{N_t-\mu_t}{1+\phi\mu_t}
}

Η ερμηνεία είναι απλή:
\begin{itemize}
  \item αν $N_t>\mu_t$, το score είναι θετικό και το επόμενο forecast
        πρέπει να κινηθεί προς τα πάνω,
  \item αν $N_t<\mu_t$, το score είναι αρνητικό και το επόμενο forecast
        πρέπει να κινηθεί προς τα κάτω,
  \item όσο μεγαλύτερη είναι η προβλεπόμενη dispersion, τόσο πιο ήπια είναι
        η διόρθωση, επειδή ένα μεγάλο error είναι λιγότερο surprising.
\end{itemize}

Το score της τρέχουσας έκδοσης είναι \textbf{normalized forecast-error
score}. Δεν είναι ο ακριβής derivative της Negative Binomial log-likelihood.
Αυτό πρέπει να παραμείνει ρητό: η φόρμουλα είναι πρακτική και σταθερή
online διόρθωση, ενώ ένα αυστηρό GAS implementation θα χρησιμοποιούσε το
scaled likelihood score της επιλεγμένης Negative Binomial parameterization.

\subsection{Η ενημέρωση state}

Τώρα η εξίσωση γράφεται με τα ονόματα των μεταβλητών του κώδικα:

\equationnote{
  s_{t+1} = \omega + \beta s_t + \alpha\,\text{score}_t
}

με την αντιστοίχιση:

\begin{center}
\begin{tabular}{lll}
\toprule
\textbf{Σύμβολο} & \textbf{Κώδικας} & \textbf{Ρόλος}\\
\midrule
$\omega$ & \code{omega} & σταθερό intercept / baseline drift\\
$\beta$ & \code{persistence} & πόσο διατηρούμε το προηγούμενο level\\
$\alpha$ & \code{learning\_rate} & πόσο αντιδρούμε στο νέο error\\
$s_t$ & \code{state} & latent log-mean του τρέχοντος forecast\\
\bottomrule
\end{tabular}
\end{center}

Η εξίσωση προκύπτει λοιπόν από τρία σχεδιαστικά desiderata:

\begin{enumerate}
  \item Θέλουμε baseline term: $\omega$.
  \item Θέλουμε μνήμη: το παλιό level $s_t$ να μην ξεχνιέται αμέσως,
        άρα $\beta s_t$.
  \item Θέλουμε correction από τα νέα δεδομένα: $\alpha\,\text{score}_t$.
\end{enumerate}

Η τρέχουσα κλάση έχει default $\beta=0.90$, $\alpha=0.10$ και $\omega=0$.
Αυτές οι τιμές είναι configuration defaults του prototype, όχι παράμετροι
που έχουν ήδη fitted από production data. Ειδικά το $\omega=0$ σημαίνει ότι
δεν έχει προστεθεί ακόμη ξεχωριστό long-run level calibration. Για production
χρήση θα πρέπει είτε να εκτιμηθεί το intercept είτε να οριστεί ρητά το
μακροχρόνιο επίπεδο που θέλουμε να διατηρεί το state.

\subsection{Αριθμητικό παράδειγμα}

Έστω ότι σε ένα bucket:

\begin{center}
  $\mu_t=100$, $N_t=120$, $\phi=0.02$, $\omega=0$,
  $\beta=0.90$, $\alpha=0.10$.
\end{center}

Τότε:

\begin{align*}
  \text{score}_t
    &= \frac{120-100}{1+0.02\cdot100}
     = \frac{20}{3}
     \approx 6.67,\\
  s_{t+1}
    &= 0 + 0.90s_t + 0.10\cdot6.67
     = 0.90s_t + 0.667,\\
  \mu_{t+1}
    &= \exp(s_{t+1}).
\end{align*}

Το νέο bucket ήταν μεγαλύτερο από το forecast, άρα η διόρθωση αυξάνει το
log-mean. Η αύξηση δεν είναι ίση με $20$ requests: περνά από normalization
και learning rate, ώστε ένα μόνο bucket να μην αλλάζει απότομα ολόκληρο το
μοντέλο.

\section{Βήμα 3: sampling του request count}

Για να δημιουργήσει ένα πιθανό μελλοντικό bucket, η μέθοδος
\code{sample\_count} χρησιμοποιεί τη Gamma--Poisson mixture representation
της Negative Binomial.

Για dispersion $\phi>0$:

\begin{align*}
  \Lambda_t &\sim \operatorname{Gamma}
    \left(k=\frac{1}{\phi},\;\theta=\phi\mu_t\right),\\
  N_t\mid\Lambda_t &\sim \operatorname{Poisson}(\Lambda_t).
\end{align*}

Το $\Lambda_t$ είναι latent Poisson rate. Η Gamma κατανομή κάνει το rate να
μεταβάλλεται από simulation σε simulation. Αυτή η επιπλέον αβεβαιότητα είναι
ακριβώς που δημιουργεί την υπερδιασπορά.

Πράγματι, για τη Gamma επιλογή:

\begin{align*}
  \mathrm{E}[\Lambda_t] &= k\theta = \mu_t,\\
  \operatorname{Var}(\Lambda_t) &= k\theta^2 = \phi\mu_t^2.
\end{align*}

Με τον νόμο total variance:

\begin{align*}
  \operatorname{Var}(N_t)
    &= \mathrm{E}[\operatorname{Var}(N_t\mid\Lambda_t)]
      +\operatorname{Var}(\mathrm{E}[N_t\mid\Lambda_t])\\
    &= \mathrm{E}[\Lambda_t]+\operatorname{Var}(\Lambda_t)\\
    &= \mu_t+\phi\mu_t^2.
\end{align*}

Όταν $\phi=0$, δεν υπάρχει Gamma mixing layer και ο κώδικας κάνει απευθείας
Poisson draw με rate $\mu_t$. Η συνάρτηση \code{\_sample\_poisson} χρησιμοποιεί
τον απλό Knuth prototype sampler για να μείνει το παράδειγμα dependency-free.

\section{Βήμα 4: fitting και online correction του rejection rate}

Η κλάση \code{RejectionRateModel} προβλέπει την πιθανότητα rejection, όχι
το workload. Από $D_t$ αποφάσεις, εκ των οποίων $R_t$ rejected, το αρχικό
estimate είναι:

\equationnote{\hat{p}_{R} = \frac{R_t}{D_t}}

Η πιθανότητα αποθηκεύεται σε log-odds state $q_t$:

\begin{align*}
  q_t &= \log\left(\frac{p_{R,t}}{1-p_{R,t}}\right),\\
  p_{R,t} &= \frac{1}{1+e^{-q_t}}.
\end{align*}

Το logistic transformation εγγυάται ότι το forecast probability βρίσκεται
πάντα στο $(0,1)$.

Μετά το κλείσιμο του bucket, το binomial forecast error είναι:

\equationnote{\text{score}^{R}_t = R_t-D_t p_{R,t}}

και η online correction είναι:

\equationnote{q_{t+1}=\beta_Rq_t+\alpha_R\text{score}^{R}_t}

Αν απορρίφθηκαν περισσότερα requests από όσα αναμένονταν, το score είναι
θετικό και αυξάνει το rejection probability. Στην compound simulation ένα
rejected request παίρνει distribution rejected, ενώ ένα approved request
παίρνει distribution approved.

\section{Βήμα 5: workload distributions χωρίς fixed request taxonomy}

Οι παρατηρήσεις του \code{workload.py} χωρίζονται μόνο με βάση το outcome:

\begin{align*}
  F_A(w) &: \text{empirical distribution workload για approved requests},\\
  F_R(w) &: \text{empirical distribution workload για rejected requests}.
\end{align*}

Δεν απαιτείται προκαθορισμένη κατηγορία τύπου request. Τα διαφορετικά
backend tasks, durations και capacity multipliers αποτυπώνονται απευθείας
στα observed weights $w_i$.

Για να έχει πρόσφατο ιστορικό μεγαλύτερη επιρροή, η παρατήρηση ηλικίας $a$
παίρνει weight:

\equationnote{\rho_a=\lambda^a, \qquad 0<\lambda\leq1}

Το $\rho_a$ είναι \textbf{στατιστικό recency weight}, όχι capacity unit.
Στο code το νεότερο observation έχει $a=0$ και weight $1$, ενώ παλαιότερα
observations έχουν weight $\lambda^a$. Η weighted empirical distribution δεν
υποθέτει Normal κατανομή και κρατά την πραγματική ασυμμετρία και τα ακριβά
tails των request weights.

Ο weighted μέσος είναι:

\equationnote{\bar{w}_A=\frac{\sum_{i\in A}\rho_iw_i}{\sum_{i\in A}\rho_i}}

και αντίστοιχα για rejected requests. Για simulation, το
\code{WeightedEmpiricalDistribution.sample} επιλέγει ένα ιστορικό observed
weight με πιθανότητες ανάλογες των $\rho_i$.

\section{Βήμα 6: η compound workload simulation}

Η \code{compound\_workload\_forecast} επαναλαμβάνει τα παρακάτω βήματα για
$M$ simulations:

\begin{enumerate}
  \item Κάνει draw request count:
        $N^{(m)}\sim\operatorname{NB}(\mu_t,\phi)$.
  \item Για κάθε request $k=1,\ldots,N^{(m)}$, κάνει outcome draw:
        $O_k^{(m)}\sim\operatorname{Bernoulli}(1-p_{R,t})$ για approval.
  \item Επιλέγει workload weight από $F_A$ αν το outcome είναι approved και
        από $F_R$ αν είναι rejected.
  \item Αθροίζει όλα τα weights:
        $W^{(m)}=\sum_{k=1}^{N^{(m)}}w_k^{(m)}$.
\end{enumerate}

Το αποτέλεσμα δεν είναι ένας αριθμός αλλά ordered sample:

\equationnote{W^{(1)},W^{(2)},\ldots,W^{(M)}}

Από αυτό μπορούμε να υπολογίσουμε:

\begin{itemize}
  \item expected workload: περίπου $\frac{1}{M}\sum_m W^{(m)}$,
  \item P50: τυπικό/median workload,
  \item P95: workload που καλύπτει περίπου το 95\% των simulated scenarios.
\end{itemize}

Το P95 είναι αυτό που χρησιμοποιεί στη συνέχεια το \code{capacity.py} ως
protected request capacity. Άρα η διαδρομή είναι:

\equationnote{
  \text{count forecast} + \text{rejection forecast} +
  \text{weight distributions}
  \;\Longrightarrow\;
  \text{distribution of total workload}
}

\section{Βήμα 7: από το workload distribution στο capacity decision}

Αν $C_t^{total}$ είναι το capacity που υπάρχει στο bucket, και
$Q_{0.95}(W_t)$ το empirical 95th percentile των simulated workloads, τότε:

\begin{align*}
  C_t^{requests} &= Q_{0.95}(W_t),\\
  C_t^{spare} &= \max(0,C_t^{total}-C_t^{requests}),\\
  C_t^{deficit} &= \max(0,C_t^{requests}-C_t^{total}).
\end{align*}

Το request capacity έχει προτεραιότητα επειδή δεσμεύεται πρώτα το
$C_t^{requests}$. Άλλα objects μπορούν να χρησιμοποιήσουν μόνο το
$C_t^{spare}$.

\section{Τι κάνει και τι δεν κάνει ακόμη ο τρέχων κώδικας}

\subsection{Τι κάνει}

\begin{itemize}
  \item Κάνει method-of-moments baseline fit για το request count.
  \item Επιτρέπει overdispersion μέσω Negative Binomial και Poisson limit.
  \item Κάνει online state correction μετά από νέο observed bucket.
  \item Μοντελοποιεί rejection ως ξεχωριστή probabilistic διαδικασία.
  \item Διατηρεί empirical request-weight distributions αντί να επιβάλλει
        Normal υπόθεση.
  \item Παράγει predictive distribution του συνολικού workload μέσω
        Monte Carlo compound simulation.
\end{itemize}

\subsection{Τι δεν κάνει ακόμη}

\begin{itemize}
  \item Η παράμετρος \code{horizon} υπάρχει στο interface, αλλά η τρέχουσα
        υλοποίηση επιστρέφει το ίδιο state για κάθε horizon. Δεν υπάρχει
        ακόμη multi-step state propagation.
  \item Το \code{score} είναι normalized error και όχι exact Negative
        Binomial likelihood score.
  \item Τα \code{persistence}, \code{learning\_rate} και \code{omega} δεν
        έχουν fitted από production history.
  \item Ο sampler Poisson είναι αναγνώσιμος prototype και όχι ο
        αποδοτικότερος sampler για πολύ μεγάλα rates.
  \item Δεν υπάρχει ακόμη queue/backlog state ή service-level model. Το
        παρόν βήμα προβλέπει demand και protected capacity.
\end{itemize}

\section{Συνοπτική αντιστοίχιση κώδικα--μαθηματικών}

\begin{center}
\begin{tabular}{p{0.35\linewidth}p{0.51\linewidth}}
\toprule
\textbf{Στοιχείο κώδικα} & \textbf{Μαθηματικός ρόλος}\\
\midrule
\code{NegativeBinomialModel.fit} & $\hat{\mu}$ και $\hat{\phi}$ από moments\\
\code{forecast} & $\mu_t=\exp(s_t)$\\
\code{update} & $s_{t+1}=\omega+\beta s_t+\alpha\text{score}_t$\\
\code{sample\_count} & Gamma--Poisson mixture\\
\code{RejectionRateModel.fit} & $\hat{p}_R=R/D$ και log-odds state\\
\code{RejectionRateModel.update} & binomial score correction\\
\code{fit\_workload\_distributions} & recency-weighted $F_A,F_R$\\
\code{compound\_workload\_forecast} & $W^{(m)}=\sum_k w_k^{(m)}$\\
\code{empirical\_quantile} & $Q_p(W)$ από simulated totals\\
\bottomrule
\end{tabular}
\end{center}

\section{Τελική εικόνα}

Ο κώδικας ακολουθεί την εξής αλυσίδα:

\begin{enumerate}
  \item Τα IAM events μετατρέπονται σε request observations.
  \item Κάθε observation μετατρέπεται σε workload weight με βάση τον ενεργό
        χρόνο, τα capacity units, τα backend tasks και το approval outcome.
  \item Τα ιστορικά counts χρησιμοποιούνται για να εκτιμηθούν το mean και η
        dispersion της ζήτησης.
  \item Το state $s_t$ μεταφέρει το forecast από bucket σε bucket και το
        score διορθώνει το state όταν έρχεται νέα πραγματικότητα.
  \item Η προσομοίωση συνδυάζει τον αριθμό requests, το rejection rate και
        τα empirical weights.
  \item Το simulated total-workload distribution μετατρέπεται σε protected
        request capacity, spare capacity και deficit.
\end{enumerate}

Η εξίσωση του $s_{t+1}$ είναι επομένως ο μηχανισμός που επιτρέπει στο count
forecast να μην είναι στατικό: διατηρεί μνήμη, βλέπει το νέο error και κάνει
ελεγχόμενη διόρθωση. Δεν υπολογίζει από μόνη της το workload. Το workload
προκύπτει μόνο στο επόμενο στάδιο, όταν το διορθωμένο count forecast
συνδυάζεται με outcomes και per-request weight distributions.

\end{document}
